% HW1 report template (236767, Winter 2026-27): Translational Motion Planning (kit 2.4).
% The report is in English and typed in LaTeX with this template. It holds, in this order:
% the AI-use box; one section for Part B, the eight planners, with one subsection per
% planner, P1 to P8; and one section for Part C, the dry question T1. Part A does not go
% into the report.
% make_submission.py submits report.pdf, or, if there is none, report_template.pdf.
%
% For each planner keep ONE of its two blocks: "flawed" (verdict, figure, principle,
% guarantee line) or "correct" (verdict), and delete the other.
% The principles and guarantee lines are GRADED FROM answers.txt (one line per key,
% see README.md); repeat them here so that the report reads on its own, and keep the
% two copies identical. The proof T1 is graded from this report.
%
% Figures come from draw.py. In the folder of this file run, for each flawed planner,
%   python draw.py --ids ID1 ID2 counterexamples/P<k>.json --planner P<k> --out figs/P<k>.pdf
% or draw all of them at once with
%   python HW1.py --ids ID1 ID2 --figs figs
% \cefig{P<k>} includes figs/P<k>.pdf (or figs/P<k>.png) when it exists, and otherwise
% prints a box with the command that makes it.
\documentclass[11pt]{article}
\usepackage[a4paper, margin=2cm]{geometry}
\usepackage[T1]{fontenc}
\usepackage{lmodern}
\usepackage{amsmath, amssymb}
\usepackage{graphicx}
\usepackage[colorlinks=true, linkcolor=blue, urlcolor=blue]{hyperref}

\setlength{\parindent}{0pt}
\setlength{\parskip}{0.4ex plus 0.3ex}
\newcommand{\code}[1]{\texttt{#1}}
\newcommand{\cefig}[1]{% \cefig{P<k>}: the figure of draw.py for planner P<k>
  \IfFileExists{figs/#1.pdf}{\includegraphics[width=\linewidth,height=8cm,keepaspectratio]{figs/#1.pdf}}{%
  \IfFileExists{figs/#1.png}{\includegraphics[width=\linewidth,height=8cm,keepaspectratio]{figs/#1.png}}{%
  \fbox{\parbox{0.95\linewidth}{\raggedright\footnotesize\texttt{figs/#1.pdf} not found; make it with\\[0.3ex]
  \texttt{python draw.py -{}-ids ID1 ID2\\ counterexamples/#1.json -{}-planner #1\\ -{}-out figs/#1.pdf}}}}}}
% \flawedblock{P<k>}{verdict}{principle}{guarantee line}: figure on the left, text on the right
\newcommand{\flawedblock}[4]{%
  \par\noindent\begin{minipage}[t]{0.48\textwidth}\vspace{0pt}\cefig{#1}\end{minipage}\hfill
  \begin{minipage}[t]{0.49\textwidth}\vspace{0pt}\raggedright
  \textbf{Verdict.} #2\par\smallskip
  \textbf{Principle.} #3\par\smallskip
  \textbf{Guarantee line.} #4
  \end{minipage}\par\medskip}
% \correctblock{note}: the verdict, with an optional one-sentence note
\newcommand{\correctblock}[1]{\par\textbf{Verdict.} Correct (no counterexample file).\par\textbf{Note.} #1\par\medskip}

\begin{document}

\begin{center}
{\large\bfseries HW1 report: Translational Motion Planning}\\[0.5ex]
Student IDs: 318906989, 323951285\qquad
Robot: $(0, 0)$, $(0, 0.6)$, $(-0.7, 0)$ % python gen_hw1.py --ids ID1 ID2 --robot
\end{center}

\noindent\fbox{\begin{minipage}{\dimexpr\textwidth-2\fboxsep-2\fboxrule}
\textbf{AI use} (at most five lines).
(i)~Tools: ChatGPT.
(ii)~Items we used them on: Understanding and analyzing the different motion planners.
(iii)~One thing a tool got wrong that we caught: While checking our guarantee for P8,
ChatGPT incorrectly claimed that overlapping C-obstacles
necessarily prevent optimality. We identified that an
optimal path may still exist if the additionally blocked
area does not affect the shortest valid path.
\end{minipage}}

\medskip
%==============================================================================
% Part B: one subsection per planner. Every subsection has two blocks; keep one.
% \flawedblock: the verdict, the figure of the counterexample, the principle the planner
% violates (one sentence: the fact from L2 that the flaw contradicts, stated as a property
% that every correct planner must have, not as a description of the bug), and its guarantee
% line (one sentence): a natural class of instances for your robot on which the planner is
% correct, described by a property of the instance alone (obstacles, start, goal), as large
% as you can make it; your own counterexample must lie outside it, and a class that contains
% any instance on which the planner is wrong earns 0. Both as in answers.txt.
% \correctblock: the verdict and an optional note; a correct planner earns its points by
% having no counterexample file.
%==============================================================================
\section*{Part B: the eight planners}

\subsection*{P1}
\flawedblock{P1}
  {Flawed. Counterexample \code{counterexamples/P1.json},
   size 1 (obstacle); \code{check\_ce.py}:
   COUNTEREXAMPLE (invalid\_path).}
  {For a translating robot, collision-free configurations must be determined by
   expanding each obstacle using the reflected robot, rather than the original
   robot.}
  {The planner is correct whenever the straight segment from start to goal does
   not pass through the interior of any true or incorrectly computed
   C-obstacle.}

\subsection*{P2}
\correctblock
  {The planner correctly computes C-obstacles, constructs the visibility graph,
   and finds a shortest path using Dijkstra's algorithm.}

\subsection*{P3}
\flawedblock{P3}
  {Flawed. Counterexample \code{counterexamples/P3.json},
   size 1 (obstacle); \code{check\_ce.py}:
   COUNTEREXAMPLE (suboptimal).}
  {A visibility graph must allow a segment along a C-obstacle boundary whenever
   the segment avoids the interior of every C-obstacle.}
  {The planner is correct when the straight path from start to goal has strictly
   positive clearance from every C-obstacle.}

\subsection*{P4}
\flawedblock{P4}
  {Flawed. Counterexample \code{counterexamples/P4.json},
   size 2 (obstacles); \code{check\_ce.py}:
   COUNTEREXAMPLE (suboptimal).}
  {A placement is forbidden exactly when it lies in the interior of at least one
   individual C-obstacle, so a shared boundary is traversable wherever it lies
   in no such interior.}
  {P4 is correct if the closed C-obstacles are pairwise disjoint, or if the
   straight start-to-goal segment has positive clearance from every closed
   C-obstacle.}

\subsection*{P5}
\flawedblock{P5}
  {Flawed. Counterexample \code{counterexamples/P5.json},
   size 0 (obstacles); \code{check\_ce.py}:
   COUNTEREXAMPLE (missed\_path).}
  {A visibility graph must include a direct edge between start and goal whenever
   the straight path is collision-free.}
  {P5 is correct if the open start-to-goal segment meets a C-obstacle interior,
   or if the segment is collision-free and contains a C-obstacle vertex strictly
   between its endpoints.}

\subsection*{P6}
\correctblock
  {The planner correctly computes C-obstacles, constructs the visibility graph
   with valid edges, and finds a shortest path using Dijkstra's algorithm with a
   priority queue.}

\subsection*{P7}
\flawedblock{P7}
  {Flawed. Counterexample \code{counterexamples/P7.json},
   size 1 (obstacle); \code{check\_ce.py}:
   COUNTEREXAMPLE (invalid\_path).}
  {A visibility graph must reject segments that pass through the interior of a
   C-obstacle, even when they intersect its boundary only at vertices.}
  {The planner is correct when the straight segment from start to goal has
   strictly positive clearance from every C-obstacle.}

\subsection*{P8}
\flawedblock{P8}
  {Flawed. Counterexample \code{counterexamples/P8.json},
   size 2 (obstacles); \code{check\_ce.py}:
   COUNTEREXAMPLE (missed\_path).}
  {A placement is forbidden exactly when it lies in the interior of at least one
   original C-obstacle, and combining obstacles must preserve this set.}
  {P8 is correct whenever the interiors of the original C-obstacles are pairwise
   disjoint, because every overlap group then contains just one obstacle.}

%==============================================================================
% Part C: the dry question T1, in at most half a page, written without relying on the
% kit's checker. answers.txt names its planner (T1.planner). Graded: 4 for the step and
% why the planner's other stages agree with the referee, 4 for why the step changes
% nothing on smaller instances, and 2 for the conclusion.
% If you skip T1, write T1.planner: - in answers.txt (T1 then earns 0) and delete this
% section's text.
%==============================================================================
\clearpage
\section*{Part C: a dry question}
\subsection*{T1: nothing smaller works}

\textbf{Planner.} P8 (as T1.planner in answers.txt).

\textbf{Proof.}
We prove that no counterexample to P8 exists with fewer than two obstacles.

P8 correctly computes each C-obstacle as the Minkowski sum of the obstacle and the reflected robot. It then constructs a visibility graph using the obstacle vertices, start, and goal, checks collision-free edges, and applies Dijkstra's algorithm to find the shortest path. By the visibility graph theorem, these steps produce an optimal path when the obstacle geometry is preserved.

The planner's only flaw is in \texttt{solid\_blocks}, which replaces overlapping C-obstacles with the convex hull of their union, potentially blocking additional free space. When this modification does not occur, the planner follows the correct shortest-path algorithm.

With zero obstacles, no merging occurs, and the straight path is optimal. With one obstacle, \texttt{solid\_blocks} returns the original C-obstacle unchanged. Therefore, the flaw cannot affect either case, and P8 is correct.

Since our counterexample contains two obstacles and demonstrates incorrect behavior, while every smaller instance is handled correctly, two obstacles are necessary and sufficient. Hence, our counterexample is minimal.
\end{document}
